\documentclass[a4paper,12pt]{article}
\usepackage{fontspec}
\usepackage{unicode-math}
\setmathfont{Latin Modern Math}
\usepackage{amsmath,amssymb}
\usepackage{graphicx}
\usepackage[hidelinks]{hyperref}
\usepackage[margin=2.5cm]{geometry}
\usepackage[numbers]{natbib}
\usepackage{booktabs,array,float}
\providecommand{\textsubscript}[1]{\ensuremath{_{\text{#1}}}}
\title{From Compute to Complexity: A Physical Theory of Intelligence Emergence}
\author{Qinrang Liu}
\date{June 2026}
\begin{document}
\maketitle
\begin{abstract}
The rapid scaling of large language models has delivered remarkable functional capabilities yet produced exponentially growing energy costs with sub-linear returns. We present the Coordination Spatiotemporal Complexity (CST) theorem, a first-principles framework that quantifies intelligence emergence through structural complexity, dynamical richness, and their physical coupling.
\end{abstract}
\section{Introduction}
Test introduction content.
\section{Results}
\subsection{The CST theorem}
We formalize the CST theorem on five axioms.
From these axioms:
\begin{equation}
CST = (S_c \cdot T_c) \cdot \exp(\alpha \cdot \Gamma_{st}) \tag{1}
\end{equation}
Spatial complexity Sc quantifies structural integration potential:
\begin{equation}
S_c = (X_1 \cdot X_2 \cdot X_3 \cdot X_4)^{1/4} \tag{2}
\end{equation}
where C is global connectivity, H is hierarchical depth, M is modularity, and R_{sw} is small-world index.
Temporal complexity Tc quantifies dynamical richness:
\begin{equation}
T_c = (\lambda_{eff} \cdot \Phi \cdot \Psi \cdot \Theta)^{1/4} \tag{3}
\end{equation}
\end{document}
